\section{One Way Door：长期正确的产品信号}

有了妙鸭和上下文设计的讨论，接下来需要一个产品筛选标准。张月光用 One Way Door 来判断一件事是否值得长期押注。

张月光喜欢 One Way Door 产品：用户一旦用了新解法，就知道自己回不去旧解法。Cursor、Claude Code 之于传统 IDE，就是典型例子。One Way Door 不一定一开始体验完美，但它让人看到未来会形成全方位优势。

\begin{figure}[H]
\centering
\includegraphics[width=0.92\textwidth]{one-way-door-product.png}
\caption{One Way Door 产品判断：旧解法、新解法、新常态。}
\end{figure}

\begin{knowledgebox}{读图：单向门不是 wow moment}
图中单向门位于旧解法和新常态之间。wow moment 只是惊艳，用户可能很快退回旧流程；单向门则改变用户的默认行为。它的市场份额可能从小开始，但会随时间持续上涨。
\end{knowledgebox}

\subsection{为什么 Agent 可能是单向门}

Agent 产品如果真的让用户突破能力边界，它就不只是提高效率，而是改变用户对“我能做什么”的理解。张月光说，Dokie 现阶段从 PPT 切入，但真正目标不是 PPT，而是帮他突破个人能力边界的 AI 朋友。

\begin{importantbox}{单向门的产品检验}
一个 AI 产品是否值得长期做，可以问：我用了以后，是否不愿回到旧方法？它是否解决旧需求，但用新方式给出全方位优势？它是否会随模型进步而越来越强？
\end{importantbox}

\subsection{本章小结}

One Way Door 是本期的产品筛选器。妙鸭是强爆款，但不一定是单向门；Agent 如果能稳定突破能力边界，则可能成为长期正确的门。

